%% lua-placeholders.tex -- Plain LuaTeX support for lua-placeholders.
%%
%% Defines the engine bridge macros the Lua core defers to, plus user-facing
%% commands that mirror the LaTeX side without requiring \newcommand,
%% \newenvironment, or any LaTeX2e package.  Load with:
%%
%%     \input lua-placeholders
%%
%% from a plain LuaTeX document.  All macros take mandatory arguments;
%% switch the active namespace with \setnamespace{name} before \param /
%% \forlistitem / etc.

% --- Defaults and bridges the Lua core renders into ------------------------
\def\curnamespace{\jobname}
\def\paramplaceholder#1{[#1]}
\def\paramnotfound#1{[<unknown> #1]}
\def\paramlistconjunction{, }

% --- Engine-specific hotfix wrappers (mirror tex/lua-placeholders.sty)
% The .sty wires these to numprint / isodate / xspace; here we keep them
% minimal -- a power user can override before \input.
\def\paramnumberformat#1{#1}
\def\paramdateformat#1#2#3{#1/#2/#3}
\def\paramfieldterm{}

% --- Boolean flag bridge: \let the if-token directly to \iftrue/\iffalse.
% \newif's internal \@if/\if@ machinery requires literal CS names and breaks
% when called with \expandafter\newif\csname if#1\endcsname, especially in
% plain LuaTeX where @ is "other" at user time.  The simpler \let approach
% works identically in plain TeX and LaTeX, supports keys with spaces, and
% has no package dependency.
\def\paramnewbool#1{%
    \expandafter\ifx\csname if#1\endcsname\relax
        \expandafter\let\csname if#1\endcsname=\iffalse
    \fi}
\def\paramsetbool#1#2{\csname paramsetbool#2\endcsname{#1}}
\def\paramsetbooltrue#1{\expandafter\let\csname if#1\endcsname=\iftrue}
\def\paramsetboolfalse#1{\expandafter\let\csname if#1\endcsname=\iffalse}

% --- User-callable macros ---------------------------------------------------
% No optional arguments (plain TeX has no \newcommand).  Set the default
% namespace via \setnamespace; per-call namespace can be passed by
% temporarily re-\def-ing \curnamespace inside a group.
\def\setnamespace#1{\def\curnamespace{#1}}
\def\strictparams{\directlua{lua_placeholders.set_strict()}}

\def\loadrecipe#1{\directlua{lua_placeholders.recipe('#1', '\curnamespace')}}
\def\loadpayload#1{\directlua{lua_placeholders.payload('#1', '\curnamespace')}}

\def\param#1{\directlua{lua_placeholders.param('#1', '\curnamespace')}}
\def\PARAM#1{\directlua{local p = lua_placeholders.param_object('#1', '\curnamespace') if p then tex.print(p:to_upper()) end}}
\def\rawparam#1{\directlua{local p = lua_placeholders.param_object('#1', '\curnamespace') if p then tex.print(p:raw_val() or p.default or p.placeholder) end}}

\def\hasparam#1#2#3{%
    \def\paramhastrue{#2}%
    \def\paramhasfalse{#3}%
    \directlua{lua_placeholders.handle_param_is_set('#1', '\curnamespace')}}

\def\ifparam#1#2#3{%
    \directlua{local p = lua_placeholders.param_object('#1','\curnamespace') if p then p:set_bool('#1') end}%
    \csname if#1\endcsname #2\else #3\fi}

\def\paramfield#1#2{\directlua{lua_placeholders.field('#1','#2','\curnamespace')}}

% paramobject: plain TeX has no \newenvironment, so the user wraps content
% with \paramobject{key}...\endparamobject (matching \begin / \end shape).
\def\paramobject#1{\directlua{lua_placeholders.with_object('#1','\curnamespace')}}
\def\endparamobject{\directlua{lua_placeholders.exit_object()}}

\def\forlistitem#1#2{\directlua{lua_placeholders.for_item('#1','\curnamespace','#2')}}
\def\fortablerow#1#2{\directlua{lua_placeholders.with_rows('#1','\curnamespace','#2')}}

% --- Load the Lua core ------------------------------------------------------
\directlua{lua_placeholders = require('lua-placeholders')}
\endinput
